\chapter{Randomness, Interaction and Proof}\label{ch:randomness}

\begin{center}\itshape
A proof can be a conversation, convincing without revealing,\\
and checkable by reading three of its bits at random.\\
\upshape(the content of this chapter, in our words)
\end{center}

In the 1970s randomness entered algorithms through number theory. Robert Solovay and Volker Strassen in 1977, and Michael Rabin in 1980, building on Gary Miller's work of 1976, gave tests that declare a number prime or composite after a few coin flips, never wrong about primes and wrong about composites with probability at most one in four per round (Tome~C2, Chapter~10). In 1977 John Gill defined the complexity classes of such algorithms. For twenty-five years primality was the showcase of randomness: a problem with a fast randomised algorithm and no known fast deterministic one, until Agrawal, Kayal and Saxena found one in 2002.

In the 1980s randomness entered \emph{proofs}. In 1985 Shafi Goldwasser, Silvio Micali and Charles Rackoff proposed to see a proof as a conversation between a prover, who may be very powerful, and a verifier, who flips coins and must not be fooled; they asked how much \emph{knowledge} a proof conveys beyond the fact that the theorem is true, and showed that some proofs convey none at all. In the same year László Babai, wanting to classify problems on matrix groups, introduced public-coin conversations between a king Arthur and a wizard Merlin. The five of them, with Shlomo Moran, shared the first Gödel Prize in 1993. Four years after their papers, Jean-Jacques Quisquater and Louis Guillou explained zero knowledge to their children, in a paper written with their families, with a cave whose two passages are joined by a door that opens only with a magic word; Peggy proves that she knows the word to Victor, who waits outside and then shouts from which passage she should come out, without ever learning it.

Then came a rush. In December 1989, in a few weeks of electronic mail circulating among a handful of researchers, Carsten Lund, Lance Fortnow, Howard Karloff and Noam Nisan showed that interactive proofs can certify that a formula is \emph{unsatisfiable}, something no short certificate is known to do, and before the month was over Adi Shamir proved that interactive proofs capture exactly polynomial space. Babai told the story in an essay with the title \emph{E-mail and the unexpected power of interaction}. Two years later, the \emph{PCP theorem} showed that every proof can be rewritten so that a verifier, flipping a few coins, reads a constant number of its bits and is convinced, with good probability, only if the proof is correct. It explained at last why so many optimisation problems resist approximation, as Tome~C2, Chapter~12, had shown by examples.

\begin{think}
You claim that two graphs drawn on paper are not isomorphic. I am a sceptic who can do polynomial computations and flip coins, and you can compute anything. Without showing me a short certificate (nobody knows one), how could you convince me?
\end{think}

\section{Randomised computation}\label{sec:bpp}

A \emph{probabilistic machine} is a deterministic polynomial-time machine $M$ that receives, besides its input $x$, a string $r$ of $p(|x|)$ random bits, uniform and independent; probabilities below are over $r$.

\begin{definition}\label{def:bpp}
A language $L$ is in
\begin{itemize}
\item $\mathsf{BPP}$ if some probabilistic machine answers correctly with probability at least $\frac23$ on every input;
\item $\mathsf{RP}$ if some probabilistic machine accepts every $x\in L$ with probability at least $\frac12$ and never accepts an $x\notin L$; $\mathsf{coRP}$ is the class of complements;
\item $\mathsf{ZPP}$ if some probabilistic machine outputs, on every input, either the correct answer or ``don't know'', the latter with probability at most $\frac12$.
\end{itemize}
\end{definition}

So $\mathsf P\subseteq\mathsf{ZPP}\subseteq\mathsf{RP}\subseteq\mathsf{BPP}$, and $\mathsf{RP}\subseteq\mathsf{NP}$ (Exercise~\ref{ex:rpnp}). Polynomial identity testing is in $\mathsf{coRP}$ by the lemma of Schwartz and Zippel (Tome~C2, Lemma~8.5): a nonzero polynomial is rarely zero at a random point, and a zero polynomial always is. In the same way Lov\'asz's test (Tome~C2, Corollary~8.7), which evaluates the determinant of the Tutte matrix of a graph at a random point, decides the existence of a perfect matching with one-sided error; it is slower than Edmonds' algorithm, but it runs in polylogarithmic parallel time, whereas the best deterministic parallel algorithms, found by Stephen Fenner, Rohit Gurjar and Thomas Thierauf (2016) for bipartite graphs and by Ola Svensson and Jakub Tarnawski (2017) in general, need a quasi-polynomial number of processors (quoted). The constants $\frac23$ and $\frac12$ do not matter.

\begin{proposition}\label{prop:amplification}
If $L\in\mathsf{BPP}$, then for every polynomial $q$ some probabilistic machine decides $L$ with error probability at most $2^{-q(n)}$ on inputs of length $n$. The same holds for $\mathsf{RP}$ with one-sided error.
\end{proposition}

\begin{proof}
Run the machine $k=48q(n)$ times with independent random strings and output the majority answer. Let $S$ be the number of correct runs; its mean is $\mu\ge\frac23k$. The majority is wrong only if $S\le\frac k2\le(1-\frac14)\mu$, which by the Chernoff bound (Tome~3, Theorem~4.6) has probability at most $e^{-\mu/32}\le e^{-k/48}=e^{-q(n)}<2^{-q(n)}$. For $\mathsf{RP}$, run the machine $q(n)$ times and accept if some run accepts: an $x\in L$ is missed with probability at most $2^{-q(n)}$, and an $x\notin L$ is never accepted.
\end{proof}

\begin{proposition}\label{prop:zpp}
$\mathsf{ZPP}=\mathsf{RP}\cap\mathsf{coRP}$.
\end{proposition}

\begin{proof}
If $L\in\mathsf{ZPP}$, answer ``yes'' when the machine says yes and ``no'' otherwise: this never accepts wrongly and accepts $x\in L$ with probability at least $\frac12$, so $L\in\mathsf{RP}$; answering ``no'' only when the machine says no shows $L\in\mathsf{coRP}$. Conversely, run an $\mathsf{RP}$ machine for $L$ and an $\mathsf{RP}$ machine for $\bar L$. If the first accepts, $x\in L$ for sure; if the second accepts, $x\notin L$ for sure; otherwise say ``don't know''. If $x\in L$, the first accepts with probability at least $\frac12$, so ``don't know'' has probability at most $\frac12$; similarly if $x\notin L$.
\end{proof}

Repeating a $\mathsf{ZPP}$ machine until it answers gives an algorithm that is always right and runs in expected polynomial time (the expected number of repetitions is at most $2$), the \emph{Las Vegas} algorithms of Tome~C2, Chapter~8.

Is randomness really useful? Many believe that $\mathsf P=\mathsf{BPP}$. Russell Impagliazzo and Avi Wigderson proved in 1997 that this follows if some problem solvable in time $2^{O(n)}$ requires circuits of size $2^{\Omega(n)}$ (quoted): hardness can be turned into pseudorandomness. Chapter~\ref{ch:circuits} returns to the converse. Meanwhile, two theorems locate $\mathsf{BPP}$ among the classes we know.

\begin{theorem}[Sipser, Gács, Lautemann 1983]\label{thm:sgl}
$\mathsf{BPP}\subseteq\Sigma_2^p\cap\Pi_2^p$.
\end{theorem}

\begin{proof}
Since $\mathsf{BPP}$ is closed under complement (exchange the answers), it is enough to show $\mathsf{BPP}\subseteq\Sigma_2^p$. Let $L\in\mathsf{BPP}$. By Proposition~\ref{prop:amplification} there is a machine $M$ using $m=m(n)$ random bits, a polynomial, with error at most $2^{-n}$. For an input $x$ of length $n$ let $A_x=\{r\in\{0,1\}^m:M(x,r)\text{ accepts}\}$. For $s\in\{0,1\}^m$ write $A_x\oplus s=\{r\oplus s:r\in A_x\}$, the set translated by the bitwise exclusive or. Let $k=\lceil m/n\rceil+1$. We claim
\[
x\in L\iff\exists s_1,\dots,s_k\in\{0,1\}^m\ \forall r\in\{0,1\}^m\ \ r\in\bigcup_{i=1}^k(A_x\oplus s_i),
\]
which is a $\Sigma_2^p$ definition, since $r\in A_x\oplus s_i$ means that $M(x,r\oplus s_i)$ accepts, checked in polynomial time for the $k$ values of $i$.

If $x\notin L$, then $|A_x|\le2^{-n}2^m$, so $k$ translates cover at most $k2^{-n}2^m<2^m$ points, as $k<2^n$ for large $n$: some $r$ is not covered. If $x\in L$, choose $s_1,\dots,s_k$ independently and uniformly. For a fixed $r$, the points $r\oplus s_i$ are independent and uniform, and $r\in A_x\oplus s_i$ if and only if $r\oplus s_i\in A_x$; so $r$ is missed by all translates with probability at most $(2^{-n})^k$. By the union bound, some $r$ is missed with probability at most $2^m2^{-nk}<1$, since $nk>m$. So some choice of $s_1,\dots,s_k$ covers everything (Figure~\ref{fig:lautemann}).
\end{proof}

\begin{figure}[ht]
\centering
\fig{ch11-lautemann}
\caption{The proof of Theorem~\ref{thm:sgl} on $\{0,1\}^8$, drawn as a $16\times16$ grid, translations being exclusive ors. Left: when the set $A_x$ of accepting random strings is large, three translates cover everything (each colour is a translate). Right: when $A_x$ is small, the same number of translates covers only a small part.}
\label{fig:lautemann}
\end{figure}

So if $\mathsf P=\mathsf{NP}$, then $\mathsf{BPP}\subseteq\Sigma_2^p=\mathsf P$ (Theorem~\ref{thm:collapse}): randomness would be useless if nondeterminism were. The second theorem replaces randomness by \emph{advice}. A language $L$ is in $\mathsf P/\mathsf{poly}$ if there are a polynomial-time machine $V$ and strings $a_0,a_1,\dots$, of length polynomial in $n$, such that $x\in L\iff V$ accepts $\langle x,a_{|x|}\rangle$. The advice $a_n$ depends only on the length of the input, and need not be computable. Equivalently (Proposition~\ref{prop:ppoly}), $L$ is decided by a family of Boolean circuits of polynomial size, one for each input length.

\begin{theorem}[Adleman 1978]\label{thm:adleman}
$\mathsf{BPP}\subseteq\mathsf P/\mathsf{poly}$.
\end{theorem}

\begin{proof}
Take a machine $M$ for $L$ with error less than $2^{-n}$ on inputs of length $n$, using $m(n)$ random bits. For each $x$ of length $n$, the set of $r$ on which $M$ errs has proportion less than $2^{-n}$. The union of these sets over the $2^n$ inputs of length $n$ has proportion less than $1$, so some $r_n$ makes $M$ correct on every input of length $n$. Take $a_n=r_n$.
\end{proof}

The proof is not constructive: finding $r_n$ requires checking all inputs of length $n$. Derandomisation is the art of finding such strings efficiently.

\section{Interactive proofs}\label{sec:ip}

An \emph{interactive proof} for a language $L$ is a protocol between a \emph{verifier} $V$, a probabilistic polynomial-time machine whose random bits are hidden from the other party, and a \emph{prover} $P$, an arbitrary function from conversations to messages. They exchange a polynomial number of messages of polynomial length, and $V$ then accepts or rejects. We require:
\begin{itemize}
\item \emph{completeness}: if $x\in L$, some prover makes $V$ accept with probability at least $\frac23$;
\item \emph{soundness}: if $x\notin L$, every prover makes $V$ accept with probability at most $\frac13$.
\end{itemize}
$\mathsf{IP}$ is the class of languages with interactive proofs. An $\mathsf{NP}$ certificate is a one-message interactive proof, so $\mathsf{NP}\subseteq\mathsf{IP}$; and $\mathsf{IP}\subseteq\mathsf{PSPACE}$, because the best acceptance probability over all provers can be computed by a recursion over the possible messages, averaging over the verifier's coins and maximising over the prover's answers, in polynomial space. As in Proposition~\ref{prop:amplification}, the constants can be improved by repetition.

Here is a problem in $\mathsf{IP}$ that is not known to be in $\mathsf{NP}$: \textup{\textsc{Graph non-isomorphism}}, the pairs $(G_1,G_2)$ of graphs on the vertex set $\{1,\dots,n\}$ that are not isomorphic.

\begin{proposition}[Goldreich, Micali, Wigderson 1986]\label{prop:gni}
\textup{\textsc{Graph non-isomorphism}} has an interactive proof with completeness $1$ and soundness $\frac14$.
\end{proposition}

\begin{proof}
Do twice, independently: the verifier picks $i\in\{1,2\}$ and a permutation $\pi$ of $\{1,\dots,n\}$, uniformly, and sends $H=\pi(G_i)$; the prover answers $j\in\{1,2\}$. The verifier accepts if $j=i$ both times.

If $G_1\not\cong G_2$, the graph $H$ is isomorphic to exactly one of them, and the prover, who can compute anything, finds which: it always answers correctly. If $G_1\cong G_2$, then $\pi(G_1)$ and $\pi(G_2)$ have the same distribution, uniform on the graphs isomorphic to $G_1$ (as $\pi$ is uniform, $\pi\circ\tau$ is uniform for a fixed isomorphism $\tau$). So $H$ carries no information about $i$, and whatever the prover does, $P(j=i)=\frac12$ in each round, $\frac14$ for both.
\end{proof}

The prover's power is used to distinguish; the verifier's secret coins are what makes cheating impossible. Can interaction certify something as hard as unsatisfiability? The key is to turn formulas into polynomials and to check a huge sum one variable at a time.

\begin{theorem}[sum-check protocol; Lund, Fortnow, Karloff, Nisan 1990]\label{thm:sumcheck}
Let $p$ be a prime, $g\in\mathbb F_p[x_1,\dots,x_v]$ a polynomial of degree at most $d$ in each variable, which the verifier can evaluate at any point in polynomial time, and $H\in\mathbb F_p$. There is an interactive proof of the claim
\[
\sum_{b\in\{0,1\}^v}g(b)=H
\]
in which the verifier accepts with probability $1$ if the claim is true, and with probability at most $vd/p$ if it is false, whatever the prover does.
\end{theorem}

\begin{proof}
The protocol has $v$ rounds (Figure~\ref{fig:sumcheck}). Set $c_0=H$. In round $i=1,\dots,v$, the prover sends a polynomial $g_i(X)$ of degree at most $d$, by its coefficients; the verifier checks that $g_i(0)+g_i(1)=c_{i-1}$, rejecting otherwise, then picks $r_i\in\mathbb F_p$ uniformly, sends it, and sets $c_i=g_i(r_i)$. At the end it accepts if $c_v=g(r_1,\dots,r_v)$. An honest prover sends
\[
s_i(X)=\sum_{b_{i+1},\dots,b_v\in\{0,1\}}g(r_1,\dots,r_{i-1},X,b_{i+1},\dots,b_v),
\]
of degree at most $d$, and then $s_1(0)+s_1(1)$ is the full sum, $s_i(0)+s_i(1)=s_{i-1}(r_{i-1})$, and $s_v(r_v)=g(r_1,\dots,r_v)$: a true claim is always accepted.

Suppose the claim is false: $c_0\ne s_1(0)+s_1(1)$. We say that round $i$ starts \emph{false} if $c_{i-1}\ne s_i(0)+s_i(1)$, writing $s_{v+1}$ for the constant $g(r_1,\dots,r_v)$ and reading the final check as round $v+1$. If round $i\le v$ starts false and the prover sends $g_i=s_i$, the check $g_i(0)+g_i(1)=c_{i-1}$ fails. Otherwise $g_i\ne s_i$; two distinct polynomials of degree at most $d$ agree on at most $d$ points, so $P\bigl(g_i(r_i)=s_i(r_i)\bigr)\le d/p$; and if $g_i(r_i)\ne s_i(r_i)$, then $c_i\ne s_i(r_i)=s_{i+1}(0)+s_{i+1}(1)$: round $i+1$ starts false. If round $v+1$ starts false, the final check fails. So the verifier accepts only if, in some round, the random $r_i$ hits one of at most $d$ bad points, which has probability at most $vd/p$ by the union bound.
\end{proof}

\begin{figure}[ht]
\centering
\fig{ch11-sumcheck}
\caption{The sum-check protocol. Each round replaces a claim about a sum over $2^{v-i+1}$ points by a claim about a sum over $2^{v-i}$ points, at a random value $r_i$; a false claim survives a round only if $r_i$ is a root of the difference between the prover's polynomial and the true one. At the end, a single evaluation of $g$ decides.}
\label{fig:sumcheck}
\end{figure}

To use it, a 3-CNF formula $\varphi$ with $n$ variables and $m$ clauses is \emph{arithmetised}: replace the literal $x_i$ by $x_i$ and $\neg x_i$ by $1-x_i$, a clause $\ell_1\lor\ell_2\lor\ell_3$ by $1-(1-\tilde\ell_1)(1-\tilde\ell_2)(1-\tilde\ell_3)$, and the conjunction by the product. The polynomial $g_\varphi$ obtained takes the value $1$ at a point of $\{0,1\}^n$ if the assignment satisfies $\varphi$, and $0$ otherwise; it has degree at most $m$ in each variable, and the verifier evaluates it anywhere in time polynomial in $m$ (Exercise~\ref{ex:arith}). So the number of models of $\varphi$ is $\sum_{b\in\{0,1\}^n}g_\varphi(b)$.

\begin{corollary}\label{cor:sharpsatip}
The language of pairs $(\varphi,K)$ such that the 3-CNF formula $\varphi$ has exactly $K$ satisfying assignments is in $\mathsf{IP}$. In particular $\mathsf{coNP}\subseteq\mathsf{IP}$.
\end{corollary}

\begin{proof}
The prover sends a prime $p$ with $3nm2^n<p\le6nm2^n$, which exists by Bertrand's postulate (Tome~1); the verifier checks its primality in polynomial time with the algorithm of Agrawal, Kayal and Saxena (Tome~C2, Chapter~10). Both $K$ and the true count lie in $[0,2^n]$, so they are equal if and only if they are equal modulo $p$. Run the sum-check protocol for $g_\varphi$ and $H=K\bmod p$: a true claim is accepted with probability $1$, a false one with probability at most $nm/p<\frac13$. Unsatisfiability is the case $K=0$.
\end{proof}

With Toda's theorem (Chapter~\ref{ch:beyondnp}) this already puts the whole polynomial hierarchy inside $\mathsf{IP}$. Shamir went all the way.

\begin{theorem}[Shamir 1990]\label{thm:shamir}
$\mathsf{IP}=\mathsf{PSPACE}$. \emph{(Quoted.)}
\end{theorem}

Shamir arithmetises a quantified formula: $\forall x$ becomes a product of the values at $0$ and $1$, $\exists x$ becomes $1-(1-\cdot)(1-\cdot)$, and the sum-check idea checks one quantifier at a time. The degrees would double with each quantifier; the trick is to insert \emph{linearisation} steps, which replace a polynomial by the one that agrees with it on $\{0,1\}^n$ and is of degree $1$ in a given variable, and which the protocol can also check. The theorem has a strange consequence: interactive proofs do not relativise, since there are oracles $A$ with $\mathsf{coNP}^A\not\subseteq\mathsf{IP}^A$ (quoted). Arithmetisation is a technique that escapes the barrier of Theorem~\ref{thm:bgs}.

\section{Zero knowledge}\label{sec:zk}

In the protocol for non-isomorphism, the verifier learns only things it knew or could compute: it gets back its own $i$. Goldwasser, Micali and Rackoff made this precise. The \emph{view} of a verifier in an interaction is its random bits together with all the messages it received; it determines everything the verifier knows at the end.

\begin{definition}\label{def:zk}
An interactive proof for $L$ with prover $P$ is \emph{perfect zero knowledge} if for every probabilistic polynomial-time verifier $V^*$, honest or not, there is a probabilistic algorithm $S$, the \emph{simulator}, running in expected polynomial time, such that for every $x\in L$ the output of $S(x)$ has exactly the distribution of the view of $V^*$ interacting with $P$ on $x$.
\end{definition}

A verifier that could produce its view alone has learnt nothing from the prover: anything it computes from the view, it could have computed from $S(x)$. The definition does not prevent learning that $x\in L$, which is the point of the proof.

\begin{theorem}[Goldreich, Micali, Wigderson 1986]\label{thm:zkgi}
\textup{\textsc{Graph isomorphism}} has a perfect zero-knowledge interactive proof with completeness $1$ and soundness $\frac12$ (and $2^{-k}$ after $k$ independent repetitions).
\end{theorem}

\begin{proof}
Let $G_1\cong G_2$ by a permutation $\tau$, known to the prover: $\tau(G_1)=G_2$. The protocol (Figure~\ref{fig:gi}):
\begin{enumerate}
\item the prover picks a uniform permutation $\pi$ and sends $H=\pi(G_1)$;
\item the verifier picks $b\in\{1,2\}$ uniformly and sends it;
\item the prover sends $\sigma=\pi$ if $b=1$, and $\sigma=\pi\circ\tau^{-1}$ if $b=2$; in both cases $\sigma(G_b)=H$;
\item the verifier accepts if $\sigma(G_b)=H$.
\end{enumerate}
\emph{Completeness} is clear. \emph{Soundness}: if $G_1\not\cong G_2$, then $H$ is isomorphic to at most one of them, so for at least one of the two values of $b$ no $\sigma$ with $\sigma(G_b)=H$ exists, and the verifier rejects with probability at least $\frac12$.

\emph{Zero knowledge.} Let $V^*$ be any verifier; its choice of $b$ may depend on $H$ and on its own random bits $\rho$. The simulator $S$ repeats: pick $\rho$ as $V^*$ would, $b'\in\{1,2\}$ and a permutation $\sigma$ uniformly, set $H=\sigma(G_{b'})$, and compute the challenge $b$ that $V^*$ would send on seeing $H$ with random bits $\rho$; if $b=b'$, output $(\rho,H,b,\sigma)$ and stop; otherwise start again.

Since $G_1\cong G_2$, the graph $H=\sigma(G_{b'})$ is uniform on the isomorphism class of $G_1$ whether $b'=1$ or $b'=2$, so $b'$ is independent of $(\rho,H)$ and $b=b'$ happens with probability exactly $\frac12$: the expected number of attempts is $2$. Conditioned on success, $(\rho,H)$ has the distribution of a real interaction ($\rho$ uniform, $H$ uniform on the class, independent), $b$ is the challenge $V^*$ computes from them, and $\sigma$ is uniform among the isomorphisms from $G_b$ to $H$, given $(\rho,H,b)$, by symmetry of the uniform choice of $\sigma$. In the real interaction, given $H$ and $b$, the prover's $\sigma$ is also uniform among the isomorphisms from $G_b$ to $H$, as $\pi$ is uniform. So the two distributions are equal.
\end{proof}

\begin{figure}[ht]
\centering
\fig{ch11-gi}
\caption{Zero knowledge for graph isomorphism. The prover knows an isomorphism $\tau$ from $G_1$ to $G_2$, and sends a random relabelling $H$ of $G_1$. Whichever graph the verifier names, the prover can show an isomorphism to $H$; each answer alone is a uniformly random permutation, and teaches nothing about $\tau$.}
\label{fig:gi}
\end{figure}

The simulator succeeds because it may \emph{rewind} the verifier, which the real prover cannot do; it plays the game of the cave of Quisquater and Guillou, filming only the runs where Peggy happened to choose the passage that Victor would name. Goldreich, Micali and Wigderson also showed that every language in $\mathsf{NP}$ has a zero-knowledge proof, if one-way functions exist: commit to a random recolouring of a $3$-colouring, and let the verifier open the two ends of one edge. That needs cryptography, the subject of Tome~C4. Zero-knowledge proofs are now used in practice, for instance to prove the correctness of computations or of transactions without revealing their inputs.

\section{Probabilistically checkable proofs}\label{sec:pcpthm}

A verifier for $\mathsf{NP}$ reads the whole certificate. What if it may read only a few bits of it, chosen at random?

\begin{definition}\label{def:pcpclass}
A language $L$ is in $\mathsf{PCP}(r(n),q(n))$ if there is a probabilistic polynomial-time verifier $V$ which, on input $x$ of length $n$ and with access to a proof string $\pi$, uses $O(r(n))$ random bits and reads $O(q(n))$ bits of $\pi$, at positions depending on $x$ and on its random bits, such that: if $x\in L$, some $\pi$ makes $V$ accept with probability $1$; if $x\notin L$, every $\pi$ makes $V$ accept with probability at most $\frac12$.
\end{definition}

With $r=\log n$ random bits, the verifier has polynomially many random strings, and can query only polynomially many positions in all; guessing the bits at those positions and checking all random strings shows $\mathsf{PCP}(\log n,1)\subseteq\mathsf{NP}$. The converse is one of the deepest theorems of complexity theory.

\begin{theorem}[the PCP theorem; Arora, Safra; Arora, Lund, Motwani, Sudan, Szegedy, 1992]\label{thm:pcpthm}
$\mathsf{NP}=\mathsf{PCP}(\log n,1)$. \emph{(Quoted.)}
\end{theorem}

Every proof of a theorem can be rewritten, with polynomial blow-up, into a form in which a constant number of random spot checks catch any error with probability $\frac12$. The original proof used algebra, low-degree polynomials and the sum-check protocol among others; in 2006 Irit Dinur gave a combinatorial proof, by repeatedly amplifying the gap of a constraint satisfaction problem with expander graphs (Tome~C1, Chapter~9), for which she received the Gödel Prize in 2019. The theorem is equivalent to a statement about approximation.

\begin{corollary}\label{cor:max3sat}
There is $\varepsilon>0$ such that, given a 3-CNF formula, distinguishing satisfiable formulas from formulas in which every assignment violates at least a fraction $\varepsilon$ of the clauses is NP-hard. Consequently, unless $\mathsf P=\mathsf{NP}$, no polynomial algorithm approximates the maximum number of satisfiable clauses within a factor $1-\varepsilon$.
\end{corollary}

\begin{proof}
Let $L$ be NP-complete with a verifier as in Theorem~\ref{thm:pcpthm}, using $c\log n$ random bits and $q$ queries, constants. Given $x$, for each of the $n^c$ random strings $\rho$, the verifier's decision is a Boolean function of the $q$ proof bits it reads; treat the proof bits as variables. A Boolean function of $q$ variables is a CNF with at most $2^q$ clauses of $q$ literals (one clause excluding each rejecting combination), and by Theorem~\ref{thm:threesat} a 3-CNF with at most $q2^q$ clauses, using new variables private to $\rho$. Let $\varphi_x$ be the conjunction of these $n^c$ groups, computed in polynomial time. If $x\in L$, some proof makes every group true, and the private variables can be set accordingly: $\varphi_x$ is satisfiable. If $x\notin L$, then for every proof at least half of the random strings reject, so for every assignment at least half of the groups contain a false clause: at least $\frac12n^c$ of the at most $q2^qn^c$ clauses are false. Take $\varepsilon=1/(q2^{q+1})$.
\end{proof}

Johan Håstad sharpened this in 2001 to the optimal constant: for formulas with exactly three literals per clause, satisfying a fraction $\frac78+\delta$ of the clauses is NP-hard for every $\delta>0$ (quoted), whereas a random assignment satisfies $\frac78$ on average (Tome~C2, Chapter~12). The hardness results for vertex cover and maximum cut quoted in that chapter, and Khot's unique games conjecture, all descend from the PCP theorem.

\section{Extracting randomness}\label{sec:extractors}

Randomised algorithms assume perfect coins. Physical sources (thermal noise, clock jitter, keystrokes) are random but biased and correlated. Can their randomness be purified? The right measure of how much randomness a source contains is the probability of its most likely value.

\begin{definition}\label{def:extractor}
A random variable $X$ with values in $\{0,1\}^n$ has \emph{min-entropy} $H_\infty(X)=\min_x\log_2\frac1{P(X=x)}$; so $H_\infty(X)\ge k$ means that no value has probability more than $2^{-k}$. The \emph{statistical distance} of two random variables $Y,Z$ with values in a finite set $\Omega$ is
\[
\Delta(Y,Z)=\max_{T\subseteq\Omega}\bigl|P(Y\in T)-P(Z\in T)\bigr|=\frac12\sum_{\omega\in\Omega}\bigl|P(Y=\omega)-P(Z=\omega)\bigr|.
\]
A function $\mathrm{Ext}:\{0,1\}^n\times\{0,1\}^d\to\{0,1\}^m$ is a \emph{strong $(k,\varepsilon)$-extractor} if, for every $X$ with $H_\infty(X)\ge k$ and $U_d$ uniform on $\{0,1\}^d$ independent of $X$, the pair $(U_d,\mathrm{Ext}(X,U_d))$ is at distance at most $\varepsilon$ from the uniform distribution on $\{0,1\}^{d+m}$.
\end{definition}

The short uniform \emph{seed} $U_d$ is necessary: no deterministic function extracts even one good bit from all sources of min-entropy $n-1$ (Exercise~\ref{ex:nodetextract}). Universal hashing (Tome~C2, Definition~3.6) gives an extractor whose seed is the choice of the hash function.

\begin{theorem}[leftover hash lemma; Impagliazzo, Levin, Luby 1989]\label{thm:lhl}
Let $\mathcal H$ be a universal family of functions $\{0,1\}^n\to\{0,1\}^m$: for $x\ne y$, $P_h\bigl(h(x)=h(y)\bigr)\le2^{-m}$. Let $X$ have min-entropy at least $k$, and $h$ be uniform in $\mathcal H$, independent of $X$. Then
\[
\Delta\bigl((h,h(X)),(h,U_m)\bigr)\le\frac12\sqrt{2^{m-k}}.
\]
In particular, if $m=k-2\log_2(1/\varepsilon)$, the map $(x,h)\mapsto h(x)$ is a strong $(k,\varepsilon)$-extractor, with seed the description of $h$.
\end{theorem}

\begin{proof}
For a random variable $D$ on a finite set $\Omega$ of size $N$, let $\mathrm{CP}(D)=\sum_\omega P(D=\omega)^2$ be its \emph{collision probability}, the probability that two independent copies of $D$ are equal. With $u$ the uniform distribution and $\delta_\omega=P(D=\omega)-\frac1N$, we have $\sum_\omega\delta_\omega=0$, hence $\sum_\omega\delta_\omega^2=\mathrm{CP}(D)-\frac1N$, and by the Cauchy--Schwarz inequality
\[
2\Delta(D,u)=\sum_\omega|\delta_\omega|\le\sqrt N\Bigl(\sum_\omega\delta_\omega^2\Bigr)^{1/2}=\sqrt{N\,\mathrm{CP}(D)-1}.
\]
Apply this to $D=(h,h(X))$ on $\Omega=\mathcal H\times\{0,1\}^m$, of size $N=|\mathcal H|2^m$. Two independent copies $(h,h(X))$ and $(h',h'(X'))$ are equal when $h=h'$, with probability $1/|\mathcal H|$, and then $h(X)=h(X')$. This happens if $X=X'$, with probability $\mathrm{CP}(X)\le2^{-k}$ (as $\sum_xP(X=x)^2\le\max_xP(X=x)$), or if $X\ne X'$ and $h$ collides on them, with probability at most $2^{-m}$ by universality. So $\mathrm{CP}(D)\le\frac1{|\mathcal H|}(2^{-k}+2^{-m})$ and $N\,\mathrm{CP}(D)-1\le2^{m-k}$, which gives the bound. With $m=k-2\log_2(1/\varepsilon)$ it is $\frac12\varepsilon\le\varepsilon$.
\end{proof}

The lemma is the reason why a password with $k$ bits of min-entropy can be turned into a key of almost $k$ uniform bits, and it is used in cryptography under the name of \emph{privacy amplification} (Tome~C4). Its seed is long, $O(n)$ bits for the family of Carter and Wegman (Tome~C2, Theorem~3.7). A random function is a strong extractor with seed length $\log_2(n-k)+2\log_2(1/\varepsilon)+O(1)$ and output close to $k+d$ bits, by a counting argument; explicit constructions approaching these parameters, such as that of Venkatesan Guruswami, Christopher Umans and Salil Vadhan (2009), use the algebra of codes (quoted). Extractors from \emph{two} independent weak sources, with no seed at all, are equivalent to explicit graphs without large cliques or independent sets: this is how Eshan Chattopadhyay and David Zuckerman obtained in 2016 the Ramsey graphs mentioned in Tome~C1, Chapter~11.

\begin{beauty}
Each section of this chapter adds a coin to a classical notion and finds it transformed. A proof checked by a sceptic who flips coins can prove unsatisfiability, and indeed anything decidable in polynomial space; it can convince without revealing; it can be checked by reading three bits. A polynomial has few roots, so a random point exposes any lie about it: that single remark, the lemma of Schwartz and Zippel, powers the sum-check protocol, the PCP theorem and much of modern cryptography. And the coins themselves need not be perfect: a little hashing turns weak randomness into almost perfect randomness.
\end{beauty}

\section{Exercises}

\begin{exercise}[$\star$]\label{ex:rpnp}
Show that $\mathsf{RP}\subseteq\mathsf{NP}$ and that $\mathsf{BPP}\subseteq\mathsf{PSPACE}$.
\end{exercise}

\begin{exercise}[$\star\star$]\label{ex:weakbpp}
Suppose that a probabilistic machine answers correctly with probability at least $\frac12+\frac1{q(n)}$, for a polynomial $q$. Show that the language it decides is in $\mathsf{BPP}$.
\end{exercise}

\begin{exercise}[$\star$]\label{ex:pnpbpp}
Show that if $\mathsf P=\mathsf{NP}$ then $\mathsf P=\mathsf{BPP}$.
\end{exercise}

\begin{exercise}[$\star\star$]\label{ex:npbpp}
Show that if $\mathsf{NP}\subseteq\mathsf{BPP}$ then $\mathsf{NP}=\mathsf{RP}$.
\end{exercise}

\begin{exercise}[$\star$]\label{ex:advice}
Show that $\mathsf P/\mathsf{poly}$ contains an undecidable language.
\end{exercise}

\begin{exercise}[$\star\star$]\label{ex:gninozk}
Show that the protocol of Proposition~\ref{prop:gni} is not known to be zero knowledge against dishonest verifiers: a verifier holding a graph $H'$ can learn from the prover whether $H'$ is isomorphic to $G_1$, to $G_2$, or to neither. How does the verifier of Theorem~\ref{thm:zkgi} differ?
\end{exercise}

\begin{exercise}[$\star$]\label{ex:arith}
Check that the arithmetisation $g_\varphi$ of a 3-CNF formula takes the value $1$ on satisfying assignments and $0$ on the other points of $\{0,1\}^n$, that its degree in each variable is at most the number of clauses containing that variable, and that it can be evaluated at a point of $\mathbb F_p^n$ with $O(m)$ arithmetic operations. Why can the prover not simply send $g_\varphi$ expanded into monomials?
\end{exercise}

\begin{exercise}[$\star$]\label{ex:pcpeasy}
Show that $\mathsf{PCP}(0,\log n)=\mathsf P$ and $\mathsf{PCP}(\log n,0)=\mathsf P$.
\end{exercise}

\begin{exercise}[$\star$]\label{ex:nodetextract}
Show that for every function $f:\{0,1\}^n\to\{0,1\}$ there is a random variable $X$ with $H_\infty(X)\ge n-1$ such that $f(X)$ is constant.
\end{exercise}

\begin{exercise}[$\star$]\label{ex:vonneumann}
A coin shows heads with an unknown probability $p\in(0,1)$, the tosses being independent. Toss it in pairs, output $0$ for (heads, tails), $1$ for (tails, heads), and nothing otherwise (von Neumann, 1951). Show that the output bits are independent and uniform, and compute the expected number of output bits per toss.
\end{exercise}

\section{Hints}

\begin{hint}{ex:rpnp}
A random string that leads to acceptance is a certificate. For $\mathsf{PSPACE}$, count.
\end{hint}
\begin{hint}{ex:weakbpp}
Repeat $O(q(n)^2)$ times and use Hoeffding's or Chernoff's inequality.
\end{hint}
\begin{hint}{ex:pnpbpp}
Theorem~\ref{thm:sgl} and Theorem~\ref{thm:collapse}.
\end{hint}
\begin{hint}{ex:npbpp}
Use the $\mathsf{BPP}$ algorithm for \textup{\textsc{Sat}} to find an assignment as in Proposition~\ref{prop:selfreduce}, then check it.
\end{hint}
\begin{hint}{ex:advice}
Unary languages: one advice bit per length.
\end{hint}
\begin{hint}{ex:gninozk}
Send $H'$ instead of a relabelling of $G_1$ or $G_2$.
\end{hint}
\begin{hint}{ex:arith}
Each clause factor is $1$ exactly when one of its literals is $1$. Count monomials.
\end{hint}
\begin{hint}{ex:pcpeasy}
Without randomness the queried positions are fixed; with $O(\log n)$ random bits and no queries, there is no proof at all.
\end{hint}
\begin{hint}{ex:nodetextract}
One of $f^{-1}(0)$, $f^{-1}(1)$ has at least $2^{n-1}$ elements.
\end{hint}
\begin{hint}{ex:vonneumann}
The two orders have the same probability $p(1-p)$.
\end{hint}

\section{Solutions}

\begin{solution}{ex:rpnp}
\idea Nondeterminism can guess the good coins; polynomial space can count them.
\writeup Let $M$ be an $\mathsf{RP}$ machine for $L$. If $x\in L$, some $r$ makes $M(x,r)$ accept; if $x\notin L$, none does. So $r$ is a certificate and $M$ the verifier: $L\in\mathsf{NP}$. For $L\in\mathsf{BPP}$, run $M(x,r)$ for all $r\in\{0,1\}^{p(n)}$ in turn, reusing the space, and count the accepting ones with a counter of $p(n)+1$ bits; accept if the count is at least half of $2^{p(n)}$. The space is polynomial.
\end{solution}

\begin{solution}{ex:weakbpp}
\idea A small advantage is amplified by many repetitions, as in a poll.
\writeup Run the machine $k=q(n)^2$ times and take the majority. Let $X_i$ be $1$ if run $i$ is correct; the $X_i$ are independent with mean at least $\frac12+\frac1q$. The majority errs only if $\sum X_i-\mathbb E\sum X_i\le-\frac kq$, which by Hoeffding's inequality (Tome~3, Theorem~4.9) has probability at most $e^{-2(k/q)^2/k}=e^{-2}<\frac13$. So the new machine, still polynomial, has error less than $\frac13$.
\end{solution}

\begin{solution}{ex:pnpbpp}
\idea Randomness lives in the second level of the hierarchy, which collapses if $\mathsf P=\mathsf{NP}$.
\writeup If $\mathsf P=\mathsf{NP}$, then $\mathsf{PH}=\mathsf P$ by Theorem~\ref{thm:collapse}, and $\mathsf{BPP}\subseteq\Sigma_2^p\subseteq\mathsf{PH}=\mathsf P$ by Theorem~\ref{thm:sgl}. The other inclusion is always true.
\end{solution}

\begin{solution}{ex:npbpp}
\idea Turn two-sided error into one-sided error by searching for a solution and checking it.
\writeup Since \textup{\textsc{Sat}} is NP-complete and $\mathsf{RP}$ is closed under polynomial reductions, it is enough to show that \textup{\textsc{Sat}} lies in $\mathsf{RP}$. Let $A$ be a $\mathsf{BPP}$ algorithm for \textup{\textsc{Sat}} with error at most $2^{-2n}$ on formulas of size $n$ (Proposition~\ref{prop:amplification}). Given $\varphi$ with $v\le n$ variables, fix the variables one by one as in Proposition~\ref{prop:selfreduce}, asking $A$ whether the formula with $x_i$ true is satisfiable; at the end, check whether the assignment found satisfies $\varphi$, and accept if and only if it does. If $\varphi$ is unsatisfiable, no assignment satisfies it, and we never accept. If $\varphi$ is satisfiable, all $v$ answers of $A$ are correct with probability at least $1-v2^{-2n}\ge\frac12$, and then the assignment found satisfies $\varphi$. So \textup{\textsc{Sat}} lies in $\mathsf{RP}$, hence $\mathsf{NP}\subseteq\mathsf{RP}\subseteq\mathsf{NP}$.
\end{solution}

\begin{solution}{ex:advice}
\idea Advice need not be computable, and one bit per length can encode an undecidable set.
\writeup Let $K\subseteq\N$ be undecidable (the codes of machines that halt on the empty input, Chapter~\ref{ch:turing}), and $U=\{1^n:n\in K\}$, undecidable too. Take as advice $a_n$ the bit ``$n\in K$'', and let $V$ accept $\langle x,a\rangle$ if $x$ is a string of $1$s and $a=1$. Then $U\in\mathsf P/\mathsf{poly}$. So $\mathsf P/\mathsf{poly}$ is not contained in the decidable languages, and $\mathsf{BPP}\subseteq\mathsf P/\mathsf{poly}$ is a statement about circuit size, not about algorithms.
\end{solution}

\begin{solution}{ex:gninozk}
\idea The prover answers questions that the verifier cannot answer itself; a dishonest verifier can ask its own questions.
\writeup Instead of a relabelling $\pi(G_i)$, the verifier sends an arbitrary graph $H'$. The honest prover answers $j$ if $H'\cong G_j$ (and something arbitrary if neither). So the verifier learns which of $G_1$, $G_2$ the graph $H'$ is isomorphic to, which it cannot compute itself in polynomial time as far as anyone knows: the protocol transmits knowledge. (Goldreich, Micali and Wigderson fix this by making the verifier first prove, in zero knowledge, that it knows $i$ and $\pi$.) In Theorem~\ref{thm:zkgi} the verifier's only message is one bit $b$, and whatever rule it uses to choose $b$, the answer it gets is a uniform isomorphism to a graph the prover chose at random, which the simulator can produce.
\end{solution}

\begin{solution}{ex:arith}
\idea Clauses become products that vanish exactly when all literals are false; evaluate the product, never expand it.
\writeup On $\{0,1\}$, $\tilde\ell$ equals the truth value of $\ell$, so $1-\tilde\ell$ is $1$ exactly when $\ell$ is false, and $1-\prod_{\ell\in C}(1-\tilde\ell)$ is $1$ exactly when some literal of $C$ is true, $0$ otherwise. The product over clauses is $1$ exactly when all clauses are satisfied. A variable occurs with degree at most $1$ in each clause factor (its literals involve distinct variables), so its degree in $g_\varphi$ is at most the number of clauses containing it. To evaluate at a point of $\mathbb F_p^n$, compute each clause factor (a few operations) and multiply the $m$ results: $O(m)$ operations modulo $p$. Expanded into monomials, $g_\varphi$ can have exponentially many terms (a product of $m$ factors with several terms each), so the verifier could not even read it; the protocol only needs evaluations at points the verifier chooses, plus univariate polynomials of degree at most $m$ from the prover.
\end{solution}

\begin{solution}{ex:pcpeasy}
\idea Without randomness or without queries, probabilistic checking collapses to deterministic computation.
\writeup In $\mathsf{PCP}(0,\log n)$ the verifier is deterministic and reads $c\log n$ positions of $\pi$; the positions may depend on bits already read, but there are at most $2^{c\log n}=n^c$ possible contents of what it reads, and we can simulate it on each of them, following the adaptive choices, accepting if some sequence of answers leads to acceptance. This is polynomial, and $x\in L$ if and only if some proof makes it accept. In $\mathsf{PCP}(\log n,0)$ the verifier reads nothing: it is a probabilistic machine with $O(\log n)$ random bits, accepting $x\in L$ with probability $1$ and $x\notin L$ with probability at most $\frac12$. Running it on all $n^{O(1)}$ random strings decides $L$ in polynomial time. Both classes contain $\mathsf P$ trivially.
\end{solution}

\begin{solution}{ex:nodetextract}
\idea A deterministic function has a big fibre, and a source can hide inside it.
\writeup One of the sets $f^{-1}(0)$, $f^{-1}(1)$ has at least $2^{n-1}$ elements; let $X$ be uniform on it. Each value has probability at most $2^{-(n-1)}$, so $H_\infty(X)\ge n-1$, and $f(X)$ is constant, at distance $\frac12$ from a uniform bit. Hence seeds, or several independent sources, are necessary.
\end{solution}

\begin{solution}{ex:vonneumann}
\idea Symmetry: heads-tails and tails-heads are equally likely, whatever the bias.
\writeup Each pair is (heads, tails) with probability $p(1-p)$, (tails, heads) with probability $(1-p)p$, and produces nothing with probability $p^2+(1-p)^2$. Given that a pair produces a bit, the bit is $0$ or $1$ with probability $\frac12$ each, and different pairs are independent, so the output bits are independent and uniform. A pair produces a bit with probability $2p(1-p)$, so the expected number of bits per toss is $p(1-p)$, at most $\frac14$ (for a fair coin). The method wastes a lot; Peter Elias showed in 1972 how to approach the entropy of the coin per toss (quoted).
\end{solution}

\section*{Sources and further reading}

\begin{itemize}
\item \emph{A research paper to read.} C.~Lund, L.~Fortnow, H.~Karloff and N.~Nisan, \emph{Algebraic methods for interactive proof systems}, Journal of the ACM 39 (1992), 859--868. Ten pages that changed complexity theory. The introduction tells how the result was found and how it surprised its authors; then comes the arithmetisation of a formula, and the protocol for the permanent of a matrix (rather than for the number of models), which uses the fact that the permanent of an $n\times n$ matrix is a sum of permanents of $(n-1)\times(n-1)$ minors, so that a claim about one permanent can be reduced to a claim about one smaller permanent by evaluating a polynomial at a random point. Compare it with our Theorem~\ref{thm:sumcheck}: the sum-check protocol is the same idea in its cleanest form, as it was extracted afterwards. The last section draws the consequences: $\mathsf{PH}\subseteq\mathsf{IP}$ via Toda's theorem, and the failure of relativisation.
\item \emph{The stories.} J.~Gill, \emph{Computational complexity of probabilistic Turing machines}, SIAM Journal on Computing 6 (1977), 675--695; S.~Goldwasser, S.~Micali and C.~Rackoff, \emph{The knowledge complexity of interactive proof systems}, SIAM Journal on Computing 18 (1989), 186--208 (STOC 1985); L.~Babai, \emph{Trading group theory for randomness}, Proceedings of the 17th ACM Symposium on Theory of Computing (1985), 421--429; L.~Babai, \emph{E-mail and the unexpected power of interaction}, Proceedings of the Fifth Annual Structure in Complexity Theory Conference (1990), 30--44; J.-J.~Quisquater, L.~C.~Guillou and others, \emph{How to explain zero-knowledge protocols to your children}, Advances in Cryptology, CRYPTO~'89, Lecture Notes in Computer Science 435 (1990), 628--631.
\item R.~Solovay and V.~Strassen, \emph{A fast Monte-Carlo test for primality}, SIAM Journal on Computing 6 (1977), 84--85; M.~O.~Rabin, \emph{Probabilistic algorithm for testing primality}, Journal of Number Theory 12 (1980), 128--138; C.~Lautemann, \emph{BPP and the polynomial hierarchy}, Information Processing Letters 17 (1983), 215--217; M.~Sipser, \emph{A complexity theoretic approach to randomness}, Proceedings of the 15th ACM Symposium on Theory of Computing (1983), 330--335; L.~Adleman, \emph{Two theorems on random polynomial time}, Proceedings of the 19th IEEE Symposium on Foundations of Computer Science (1978), 75--83; R.~Impagliazzo and A.~Wigderson, \emph{P = BPP if E requires exponential circuits: derandomizing the XOR lemma}, Proceedings of the 29th ACM Symposium on Theory of Computing (1997), 220--229.
\item A.~Shamir, \emph{IP = PSPACE}, Journal of the ACM 39 (1992), 869--877; O.~Goldreich, S.~Micali and A.~Wigderson, \emph{Proofs that yield nothing but their validity, or all languages in NP have zero-knowledge proof systems}, Journal of the ACM 38 (1991), 691--729; S.~Arora and S.~Safra, \emph{Probabilistic checking of proofs: a new characterization of NP}, Journal of the ACM 45 (1998), 70--122; S.~Arora, C.~Lund, R.~Motwani, M.~Sudan and M.~Szegedy, \emph{Proof verification and the hardness of approximation problems}, Journal of the ACM 45 (1998), 501--555; I.~Dinur, \emph{The PCP theorem by gap amplification}, Journal of the ACM 54 (2007), article 12; J.~Håstad, \emph{Some optimal inapproximability results}, Journal of the ACM 48 (2001), 798--859.
\item R.~Impagliazzo, L.~A.~Levin and M.~Luby, \emph{Pseudo-random generation from one-way functions}, Proceedings of the 21st ACM Symposium on Theory of Computing (1989), 12--24; J.~von Neumann, \emph{Various techniques used in connection with random digits}, National Bureau of Standards Applied Mathematics Series 12 (1951), 36--38; V.~Guruswami, C.~Umans and S.~Vadhan, \emph{Unbalanced expanders and randomness extractors from Parvaresh--Vardy codes}, Journal of the ACM 56 (2009), article 20; E.~Chattopadhyay and D.~Zuckerman, \emph{Explicit two-source extractors and resilient functions}, Annals of Mathematics 189 (2019), 653--705; S.~Vadhan, \emph{Pseudorandomness}, Foundations and Trends in Theoretical Computer Science 7 (2012), 1--336.
\item S.~Arora and B.~Barak, \emph{Computational Complexity: A Modern Approach}, Cambridge University Press, 2009, Chapters~7, 8, 11 and~21; O.~Goldreich, \emph{Computational Complexity: A Conceptual Perspective}, Cambridge University Press, 2008, Chapters~6 and~9.
\end{itemize}
